% ---------------------------------------------------------------
% Section 4 — Implementation in the BX3 Framework
% Paper: Bailout Protocol (Pillar 5)
% ---------------------------------------------------------------

\section{Implementation in the BX3 Framework}
\label{sec:bailout-implementation}

The Bailout Protocol is not a safety module appended to an existing agent runtime.
It is the \emph{runtime expression} of the BX3 Framework's upstream accountability
guarantee --- the mechanism that converts a declarative layer commitment into a binding,
non-configurable, auditable operational decision at the precise moment a decision
proposal crosses a defined boundary.  This section details how the protocol integrates
with the Purpose Layer, how the Bounds Engine triggers bail-out events, how the Fact Layer
enforces the resulting audit trail, how the state machine governs the full trigger
lifecycle, and the formal relationship between the protocol and BX3's upstream intent.

% ----------------------------------------------------------------
\subsection{Purpose Layer as Human Accountability Anchor}
% ----------------------------------------------------------------

The Purpose Layer (\texttt{PurposeLayer.ts}) defines the intent, judgment, and
accountability boundary for every node in the BX3 execution tree.  Each node is
initialised with three canonical, immutable fields:

\begin{itemize}
  \item \texttt{nodeId} --- a unique identifier for this agent or process.
  \item \texttt{humanRoot} --- the designated human decision-maker at the terminal
        boundary of the accountability chain for this node's subtree.
  \item \texttt{accountabilityAnchor} --- a stable, versioned reference string that
        identifies the governing policy regime under which this node operates.
\end{itemize}

These fields are sealed at node initialisation and propagated downward through the
entire execution tree.  When a child node is spawned, it receives a \texttt{parentId}
and a \texttt{parentPurposeHash} that cryptographically binds it to its parent's
purpose.  The chain is not merely informational --- it is tamper-evident: any divergence
between a child's purpose and its parent's purpose is detectable by comparing the
stored hash against the parent's computed hash at any point in the execution history.

The Bailout Protocol leverages the Purpose Layer at two distinct stages:

\begin{enumerate}
  \item \textbf{Anchor resolution at trigger time.}  When the protocol fires, it
        traverses the \texttt{propagationChain} upward until it locates the node
        whose \texttt{nodeId} matches the \texttt{humanRoot} field.  This node is the
        mandatory review endpoint.  No autonomous resolution is permitted beyond this
        boundary.  The anchor resolution is deterministic and $O(1)$ relative to tree depth.

  \item \textbf{Provenance recording at fire time.}  The current
        \texttt{accountabilityAnchor} is written into the trigger's \texttt{event}
        payload at the moment of firing.  This permanently attaches the governing
        policy context to the bail-out record, enabling post-hoc reconstruction of
        which policy regime was active at the time of the event --- independent of
        whether that policy has since been revised.  A policy change after the fact
        cannot retroactively alter the anchor that was recorded.
\end{enumerate}

The Purpose Layer thus provides the \emph{semantic grounding} for every bail-out.
It answers not only ``who is responsible?'' but also ``under which governing
commitment was this escalation mandatory?''  This distinction is essential for
regulatory defensibility: an auditor can reconstruct the full policy context
surrounding any trigger, not merely its operational chain.

% ----------------------------------------------------------------
\subsection{Bounds Engine: The Triggering Mechanism}
% ----------------------------------------------------------------

The Bounds Engine (\texttt{BoundsEngine.ts}) implements BX3's bounded-reasoning
layer.  Its responsibilities are two-fold: (i) generating decision proposals
(\emph{P5}), and (ii) projecting the consequences of those proposals through
confidence scoring (\emph{P6}).  The Bounds Engine is the exclusive source of
bail-out trigger events in the BX3 runtime.  No other component may directly invoke
\texttt{BailoutProtocol.fire()}.  This exclusivity is architecturally enforced:
the engine is the only module with a dependency on the Bailout Protocol's internal API.

Triggering occurs when the Bounds Engine evaluates a proposal against the active
\texttt{SafetyEnvelope} and detects one of three named \texttt{TriggerCondition}
types:

\begin{verbatim}
type TriggerCondition =
  | 'capability_boundary'
      // proposed action exceeds the defined capability of this node
  | 'safety_envelope_predicted'
      // P6 projection forecasts a SafetyEnvelope violation
  | 'accountability_boundary'
      // proposed action requires human decision by active policy
\end{verbatim}

Each condition is evaluated against the \texttt{SafetyEnvelope} struct, which carries
domain-specific constraints (e.g., water-rights limits, soil-moisture thresholds,
equipment load limits).  Critically, the Bounds Engine does not block actions directly ---
it generates a \texttt{BailoutTrigger} record and returns control to the Bailout Protocol
orchestrator.  This separation is architectural: the Bounds Engine reasons about
possibility; the Bailout Protocol governs what happens when possibility exceeds permission.

When a trigger condition is met, the protocol calls \texttt{BailoutProtocol.fire()}:

\begin{verbatim}
fire({
  triggerCondition: TriggerCondition,
  confidence: number,           // P6 confidence score
  event: Record<string,unknown>,
  nodeState: Record<string,unknown>,
  proposedResolution: string | null
}): BailoutTrigger
\end{verbatim}

The returned \texttt{BailoutTrigger} carries a unique \texttt{triggerId}, the
\texttt{originatingNodeId}, the full \texttt{propagationChain}, and the
\texttt{accountabilityAnchor} from the Purpose Layer at fire time.  The trigger
is now a first-class runtime object managed by the state machine described in
Section~\ref{sec:state-machine}.

% ----------------------------------------------------------------
\subsection{Fact Layer: Enforcement Partner and Audit Authority}
% ----------------------------------------------------------------

The Fact Layer (\texttt{FactLayer.ts}) serves two roles in the Bailout Protocol:
the \emph{enforcement partner} that gates actionuation, and the \emph{audit authority}
that maintains the forensic ledger.

\bigskip
\noindent\textbf{Fact Layer as enforcement partner.}  The Fact Layer's
\texttt{evaluate(P5Decision, P6Projection): SandboxVerdict} method is the
architectural choke point for all P5-to-P8 transitions.  It accepts a
\texttt{P5Decision} and its corresponding P6 projection, then returns one of
three verdicts:

\begin{verbatim}
type SandboxVerdict = 'EXECUTE' | 'BLOCK' | 'REVIEW'
\end{verbatim}

The \texttt{EXECUTE} verdict is a necessary but not sufficient condition for
actionuation.  Even after a proposal passes the Sandbox Gate, the Bailout Protocol
may suppress execution if the trigger is in \texttt{human\_review} status.
The Fact Layer is therefore not the sole authority --- it is the enforcement partner
of a two-key system: Fact Layer clears the technical path; the Bailout Protocol
clears the accountability path.  Both keys must be turned.

\bigskip
\noindent\textbf{Forensic ledger.}  The Fact Layer maintains a tamper-evident
forensic ledger (\texttt{ForensicLedger}) that records every state transition,
trigger event, and resolution action with cryptographic integrity seals.
The ledger's append-only structure is:

\begin{verbatim}
seal_n = SHA256(
  JSON.stringify(payload_n) ||
  previousSeal_{n-1} ||
  timestamp_n
)
\end{verbatim}

where \texttt{previousSeal\_\{n-1\}} is the \texttt{seal} of the entry immediately
preceding the new one (or the string literal \texttt{'GENESIS'} for the first entry).
This construction is equivalent to a blockchain block header: any modification of a
historical entry invalidates all subsequent seals, because each subsequent entry's
computed \texttt{seal} would not match the stored value.

For the Bailout Protocol specifically, the ledger records four event types:

\begin{enumerate}
  \item \texttt{TRIGGER\_FIRED} --- logged at P1 when \texttt{fire()} is called,
        capturing the full \texttt{triggerCondition}, \texttt{confidence}, and
        \texttt{nodeState} snapshot.  This entry anchors the entire trigger record
        in time.

  \item \texttt{ESCALATED} --- logged at each P3--P5 hop, recording the intermediate
        \texttt{nodeId}, the action taken (\texttt{ESCALATED}), and the
        \texttt{propagationDirection}.  The sequence of \texttt{ESCALATED} entries
        reconstructs the exact path the trigger took toward the \texttt{humanRoot}.

  \item \texttt{HUMAN\_REVIEW} --- logged at P7 when the trigger arrives at the
        \texttt{humanRoot}, recording the arrival timestamp, the trigger's
        \texttt{propagationChain} length at arrival, and a \texttt{pendingDecision}
        field set to \texttt{true}.

  \item \texttt{RESOLUTION\_RECORDED} --- logged at P8 when the human resolves the
        trigger, recording the \texttt{resolutionText}, the resolver's identity
        (derived from \texttt{humanRoot}), and the end-to-end latency computed as
        $\delta t = \texttt{resolvedAt} - \texttt{firedAt}$.  This entry closes the
        trigger record and establishes the total time the system was in a
        constrained state.
\end{enumerate}

The combination of the \texttt{propagationChain} array inside the \texttt{BailoutTrigger}
and the corresponding \texttt{LedgerEntry} sequence in the Forensic Ledger provides
dual-mode auditability: the chain gives a human-readable provenance graph, while the
ledger gives a machine-verifiable cryptographic proof of integrity.  Both are generated
from the same runtime events, eliminating the possibility of divergence between
operational records and forensic records.

% ----------------------------------------------------------------
\subsection{State Machine as Fact Layer Extension}
\label{sec:state-machine}
% ----------------------------------------------------------------

The Bailout Protocol is implemented as a deterministic state machine whose states
form a superset of the BX3 layer-plane taxonomy.  Five states are defined, each
mapped to a \texttt{status} field in the \texttt{BailoutTrigger} record:

\bigskip
\noindent
\begin{tabular}{llp{7.5cm}}
  \hline
  \textbf{State} & \textbf{Plane} & \textbf{Semantics} \\
  \hline
  \texttt{pending}         & P1 & Trigger authored; \texttt{propagationChain} initialised with
                              originating node; not yet transmitted \\
  \texttt{escalating}    & P3--P5 & Propagating upward through intermediate nodes;
                              each hop appends to \texttt{propagationChain} \\
  \texttt{human\_review} & P7 & Arrived at \texttt{humanRoot}; execution is blocked;
                              human decision is required and pending \\
  \texttt{resolved}      & P8 & Human has acted; resolution recorded; trigger closed \\
  \texttt{expired}       & P9 & TTL exceeded without human resolution; protocol
                              enters a defined failure state \\
  \hline
\end{tabular}
\bigskip

State transitions are driven by three primitive operations:

\begin{itemize}
  \item \texttt{fire()} --- P1 transition.  Creates the \texttt{BailoutTrigger} record,
        initialises \texttt{propagationChain} with the originating node, sets
        \texttt{status} to \texttt{pending}, and appends a \texttt{TRIGGER\_FIRED}
        entry to the Forensic Ledger.

  \item \texttt{receiveAndPropagate()} --- P3--P5 transitions.  Each intermediate node
        in the propagation chain receives the trigger, appends its \texttt{nodeId},
        action (\texttt{ESCALATED}), and timestamp to \texttt{propagationChain},
        advances \texttt{status} to \texttt{escalating}, and forwards the trigger
        toward the next node in the chain.  A corresponding \texttt{ESCALATED}
        ledger entry is written at each hop.

  \item \texttt{resolve(\{resolutionText, resolverId\})} --- P7/P8 transition.
        The \texttt{humanRoot} node receives the trigger, changes \texttt{status}
        to \texttt{human\_review} (P7), records the human's decision, and advances
        \texttt{status} to \texttt{resolved} (P8).  A \texttt{RESOLUTION\_RECORDED}
        entry is appended to the Forensic Ledger, closing the trigger record.
\end{itemize}

The state machine is a Fact Layer extension in the specific sense that all of its
state variables --- \texttt{status}, \texttt{propagationChain}, \texttt{firedAt},
\texttt{resolvedAt} --- are maintained within the Fact Layer's authority domain.
No Bounds Engine or Purpose Layer process can write or mutate these fields.
The state machine's integrity is therefore guaranteed by the same deterministic
enforcement mechanism that guarantees all Fact Layer state transitions.

% ----------------------------------------------------------------
\subsection{Relationship: Bailout Protocol as Runtime Expression of
            the Upstream Accountability Guarantee}
% ----------------------------------------------------------------

The BX3 Framework's accountability guarantee is formally expressed in the Intent
Layer's constraint: \emph{for every intent $i \in I$, there exists a role $r \in R$
and an action $a \in A$ such that $i \rightarrow (r, a)$}.  This guarantee is
declarative --- it describes what must be true about the system's design.  The Bailout
Protocol is the mechanism by which this declaration becomes operative at runtime.

The causal chain from declaration to enforcement is:

\begin{enumerate}
  \item The \texttt{humanRoot} field in the Purpose Layer \emph{declares} that a
        specific human is the final accountable authority for actions originating
        under this node's governance subtree.
  \item The \texttt{accountability\_boundary} condition in the Bounds Engine
        \emph{detects} when a proposed action falls outside the scope of autonomous
        authorisation --- i.e., when the action is required by policy to reach the human.
  \item The Bailout Protocol's \texttt{fire()} method \emph{operationalises} the
        detection by creating a trigger that is obligated to propagate to the
        declared human.  The trigger's propagation is not discretionary; the state
        machine enforces it.
  \item The state machine \emph{enforces} that no autonomous resolution is
        permitted beyond \texttt{human\_review}.  Even if the Fact Layer has cleared
        the technical path, the accountability path blocks execution until the human
        resolves.
  \item The Forensic Ledger \emph{records} every step of this chain, providing
        cryptographic evidence that the guarantee was honoured in practice.
\end{enumerate}

If any step in this chain is absent --- if the Bounds Engine fails to detect the
condition, if the protocol fails to propagate, if the state machine permits
autonomous override, or if the ledger is not written --- the upstream accountability
guarantee is voided.  The Bailout Protocol is therefore not an additional feature
layered on top of BX3; it is the \emph{necessary and sufficient runtime
implementation} of the guarantee.  Without it, the BX3 Intent Layer is a
specification without an execution path.

This relationship is the basis for regulatory defensibility.  An auditor who accepts
the BX3 Intent Layer's formal guarantee as a design commitment can verify its
enforcement by querying the Forensic Ledger for any trigger whose
\texttt{triggerCondition} is \texttt{accountability\_boundary} and confirming that
the \texttt{propagationChain} terminates at the declared \texttt{humanRoot} with a
corresponding \texttt{RESOLUTION\_RECORDED} entry.  The existence of a queryable,
cryptographically integrity-verified record is itself evidence of the guarantee's
operationalisation --- not merely its declaration.

The protocol thus closes the loop between BX3's abstract layer model and its
operational reality: the same formal structure that defines what each layer does
also defines how accountability is enforced when layers interact in ways that
violate declared constraints.  The Bailout Protocol is the point of translation
between BX3's architectural commitments and their runtime consequences.

% ---------------------------------------------------------------
% End of Section 4
% ---------------------------------------------------------------